\section{Actividades, precedencias y cargas de innovación}
Las actividades se derivan de los paquetes del T-14 y conservan su código de cuenta. La Tabla~\ref{tab:sd7-innov-curva} agrega las 2.352 HH de mes fijo, con los períodos de origen de SD13; la Tabla~\ref{tab:sd7-innov-sobrecargas} contrasta esas cargas con migración y ensayos/implantación ya estimados. El piloto estacional de IN2, las revisiones por embarcación y los servicios recurrentes se agregan al fijar sus meses efectivos; no se consideran horas libres.

Cada aporte estimado por perfil se deriva en una actividad identificada como código del paquete más \texttt{ACT-perfil}. Funcional define o valida reglas; Datos prepara y comprueba cobertura; Arquitectura/Integración especifica o verifica interfaces; Desarrollo implementa, prueba unidades y documenta; Seguridad especifica o verifica sus controles; Implantación prepara usuarios y registra el piloto; Operación comprueba continuidad y transferencia; Proyecto controla decisiones contractuales; y Calidad consolida comprobaciones y evidencia. Cada actividad conserva su ejecutor, el responsable último de la ficha y las horas ya estimadas, sin añadir porcentajes de gestión.

Las actividades dependen de las predecesoras de la ficha; la consolidación de Calidad depende además de los aportes técnicos y funcionales del mismo paquete. Podrán coordinarse en paralelo cuando existan insumos disponibles; las dependencias entre paquetes, esperas y dedicaciones se incorporarán a la red completa. El registro no presupone que todas las personas estén disponibles simultáneamente.

El registro asociado \path{componentes/actividades_innovaciones.csv} contiene las actividades, sus horas por perfil, lugar, período y predecesoras. La construcción se realizará en DEV/QA; sus pruebas específicas, en QA/PREPROD; las observaciones del piloto y capacitación, con usuarios de la marina en ventana autorizada, y el análisis y expediente, en los entornos de Synaptix. La asignación de una quincena es una ventana de control: si la disponibilidad efectiva no permite terminar el resultado dentro de ella, se reprogramarán actividades o se descompondrá la ficha antes de fijar la línea base.

Las dependencias BASE-CTR, BASE-VAR, BASE-NAV y BASE-AMB significan contratos, modelo y controles de Clientes y contratos, Varadero, Operación marítima y Ambiente y consumos, respectivamente. En diseño exigen especificaciones aprobadas; para integración y piloto, servicios y datos validados. No se interpretan como módulos aceptados por escribir su código en la ficha. IN2 depende de IN1 habilitada; IN4 comparte el registro de tiempos con IN1, con indicadores y evidencia diferenciados; IN5 requiere medición individual y permanencias conciliadas. La coordinación de la captura conserva universos distintos y no vuelve a sumar la misma tarea.

\begingroup\small\setlength{\tabcolsep}{2pt}
\begin{longtable}{R{9mm}R{12mm}R{12mm}R{12mm}R{12mm}R{12mm}R{12mm}R{14mm}R{12mm}R{14mm}}
\caption{Cargas mensuales incrementales de innovación con mes fijo}\label{tab:sd7-innov-curva}\\\hline
Mes & A & D & S & Q & V & I & O+L & J & HH\\\hline\endfirsthead
\hline Mes & A & D & S & Q & V & I & O+L & J & HH\\\hline\endhead
2 & 16 & 0 & 0 & 8 & 0 & 0 & 0 & 0 & 24 \\\hline
4 & 16 & 0 & 0 & 8 & 0 & 0 & 0 & 8 & 32 \\\hline
14 & 64 & 16 & 24 & 24 & 8 & 56 & 0 & 0 & 192 \\\hline
15 & 16 & 16 & 40 & 16 & 0 & 48 & 0 & 0 & 136 \\\hline
16 & 0 & 32 & 0 & 56 & 272 & 8 & 0 & 0 & 368 \\\hline
17 & 32 & 0 & 40 & 96 & 384 & 48 & 0 & 0 & 600 \\\hline
18 & 24 & 24 & 56 & 160 & 56 & 8 & 48 & 0 & 376 \\\hline
19 & 48 & 24 & 0 & 88 & 0 & 0 & 32 & 0 & 192 \\\hline
20 & 64 & 40 & 16 & 136 & 0 & 0 & 40 & 16 & 312 \\\hline
21 & 8 & 0 & 0 & 24 & 0 & 0 & 80 & 8 & 120 \\\hline
\end{longtable}\FuenteTabla{Agregación de fichas de innovación con mes fijo; excluye piloto estacional, revisión por embarcación y operación recurrente. O+L comparten persona.}\endgroup

\begin{table}[htbp]\centering
\begin{tabularx}{\linewidth}{R{12mm}L{30mm}R{22mm}R{22mm}Y}\hline
Mes & Perfil & HH conocidas & Exceso & Capacidad de contraste\\\hline
16 & V & 272 & 128 & 144 HH/persona/mes \\\hline
17 & Q & 156 & 12 & 144 HH/persona/mes \\\hline
17 & V & 384 & 240 & 144 HH/persona/mes \\\hline
18 & Q & 168 & 24 & 144 HH/persona/mes \\\hline
20 & Q & 152 & 8 & 144 HH/persona/mes \\\hline
21 & O+L & 164 & 20 & 144 HH/persona/mes \\\hline
\end{tabularx}\caption{Sobrecargas detectadas antes de incorporar las restantes cuentas}\label{tab:sd7-innov-sobrecargas}
\FuenteTabla{Suma por persona de migración, ensayos/implantación, innovaciones con mes fijo y su seguimiento recurrente fijo. Excluye aún IN2 estacional y demás cuentas. Hipótesis de capacidad productiva mensual del T-15.}\end{table}


El calendario de origen concentra Desarrollo en 16--17 y excede la capacidad de una persona incluso antes de agregar el software obligatorio E2. Calidad también presenta excesos en 17, 18 y 20. En el 21, la coincidencia de transferencia final, estabilización de datos y seguimiento fijo concentra 164 HH en José Lara Arce frente a 144, en este contraste parcial que excluye IN2 variable y el servicio general; sus paquetes fijos sí se incorporan en la asignación conjunta posterior. Por ello, estas fechas son entradas para nivelación, no una línea base aprobada. La consolidación fija posterior propone refuerzos y comprueba los enlaces resueltos de esfuerzo directo. La Dirección completará las puertas, las esperas, la provisión y el calendario, y reconciliará cualquier anticipo con SD13. Si esa combinación no alcanza, deberá dimensionar recursos adicionales y su costo antes de comprometer la línea base. Ninguna corrección podrá acortar pilotos, eliminar controles, duplicar personas con funciones compartidas o desplazar hitos contractuales \parencite[art.~17; Formulario~T-15; pp.~12--63]{bases-admin}.

Las 2.352 HH de innovación con mes fijo se suman a las 2.484 HH de implementación de migración y ensayos/implantación: subtotal 4.836 HH. Operación reúne 96 de estabilización de datos, 328 de piloto estacional y 2.304 de seguimiento incremental fijo: subtotal 2.728 HH, más $16(57-M_{\mathrm{inv}})$ de IN2, $16\sum B_{o,q}$ de revisión de botaduras y la cobertura general adicional. Estos subtotales parciales excluyen revisión documental variable y otras cuentas; el balance conjunto incorpora los 4.832 HH fijos de gobierno y servicio operacional y alcanza 7.560 HH de operación, sin absorber sus variables ni guardias. La habilitación final de innovaciones del mes 21 permanece dentro de implementación; su servicio recurrente comienza en operación y no se cuenta como otra habilitación.

La curva recurrente fija de innovación aporta, en cada mes 21--56, Funcional 12, Datos 12, Calidad 16, Operación 20 y Proyecto 4 HH: 64 por mes y 2.304 en los 36 períodos. IN2 añadirá desde $M_{\mathrm{inv}}$ Funcional 4, Calidad 4 y Operación 8 HH por mes, total 16. Las horas del piloto y revisión de botaduras se agregan sólo en sus quincenas efectivas. El registro \path{componentes/servicio_innovaciones.csv} instancia los 216 paquetes fijos de seguimiento; IN2 conserva su fórmula hasta fijar el calendario. Son cargas de resultados incrementales, sin cobertura implícita de guardias o atención continua.
